\documentclass[11pt,a4paper]{article}
\usepackage[utf8]{inputenc}
\usepackage[T1]{fontenc}
\usepackage{amsmath,amssymb}
\usepackage{booktabs}
\usepackage{listings}
\usepackage{xcolor}
\usepackage{hyperref}
\usepackage{geometry}
\geometry{margin=2.5cm}

\lstset{
  basicstyle=\ttfamily\small,
  breaklines=true,
  frame=single,
  backgroundcolor=\color{gray!10},
  numbers=left,
  numberstyle=\tiny
}

\title{02244 - Logic for Security\\
\vspace{0.5em}
Project on Security Protocols}

\author{%
  \textbf{Authors}\\[0.5em]
  Zeinab\\
  Fahmida\\
  Mirtha\\[1em]
  March 2026
}

\date{}

\begin{document}
\maketitle
\thispagestyle{empty}

\newpage
\tableofcontents
\newpage

\section{Introduction - Zeinab}

\noindent
This project develops an Open Authorization protocol for the following setting.
User $A$ stores photos at server $B$ and wants to use a photo-book service $P$.
Instead of downloading from $B$ and uploading to $P$, $A$ should authorize $P$ to access selected photos directly at $B$ with limited rights.

The trust architecture includes an identity provider (\textit{IDP}, denoted $I$).
User $A$ authenticates through a password-based relation with $I$, and authorization tokens are issued and signed by $I$.
The development follows the weekly assignment structure:
Dolev--Yao baseline analysis, lazy intruder executability, typed formats and key lookup, pseudonymous channels, and guessable-password verification.

This report presents design choices, formal AnB models, and OFMC verification outcomes.
Intermediate vulnerable versions are intentionally included because they justify each refinement step.

\newpage
\section{Open Authorization Main Protocol - Zeinab}

\subsection{Outline of the design}

\subsubsection{Agents and initial knowledge}
The main agents are:
\begin{itemize}
  \item \textbf{$A$ (Owner):} delegates access.
  \item \textbf{$B$ (Photo Server):} serves protected photos.
  \item \textbf{$P$ (Photo Service):} requests delegated read access.
  \item \textbf{$I$ (IDP):} trusted authorization issuer.
\end{itemize}

The IDP is assumed honest.
Password material is used only for authentication evidence and not as an encryption key.
Authorization context includes identities and scope metadata.

\subsubsection{Security properties to establish}
\begin{itemize}
  \item secrecy of photo payload between $B$ and $P$,
  \item authentication of server decisions based on IDP-issued grants,
  \item robustness against guessable-password attacks (Week 6 task).
\end{itemize}

\subsection{AnB formalization}

\subsubsection{Design decisions and development}
Week 2 starts with a minimal baseline.
After static Dolev--Yao analysis exposed leakage risk, Week 3 introduces IDP-issued grants.
The revised flow is:
\begin{enumerate}
  \item $P \rightarrow A$: request with scope and freshness.
  \item $A \rightarrow I$: request plus password witness $h(pw(A,I),N)$.
  \item $I \rightarrow A$: signed grant.
  \item $A \rightarrow P$: token forwarding.
  \item $P \rightarrow B$: token presentation.
  \item $B \rightarrow P$: encrypted data return.
\end{enumerate}

\subsubsection{Modeling in AnB}
\begin{lstlisting}
Protocol: OpenAuth_W3_DelegationViaIDP

Types:
  Agent A,B,P,I;
  Number F,T,NA,Grant,Photos;
  Function pk,pw,h;

Actions:
  P -> A: P,B,F,T,NA
  A -> I: A,P,B,F,T,NA,h(pw(A,I),NA)
  I -> A: {A,P,B,F,T,NA,Grant}inv(pk(I))
  A -> P: {A,P,B,F,T,NA,Grant}inv(pk(I))
  P -> B: {A,P,B,F,T,NA,Grant}inv(pk(I)),F
  B -> P: {Photos}pk(P)
\end{lstlisting}

\newpage
\section{Formats and GetPublicKey Protocol - Fahmida}

\subsection{Outline of the design}

\subsubsection{Agents and initial knowledge}
Week 4 keeps the same trust model but improves message typing.
Distinct format constructors are used to separate protocol intents.
Additionally, $A$ obtains unknown public keys from $I$ via a dedicated key lookup protocol.

\subsubsection{Security properties to establish}
\begin{itemize}
  \item reduce type-flaw risk through constructor separation,
  \item authenticate key responses from $I$ to $A$.
\end{itemize}

\subsection{AnB formalization}

\subsubsection{Design decisions and development}
Heavily wrapped format variants initially produced non-executability in OFMC.
The model was simplified while preserving typed message structure.
A separate textual argument documents type-flaw resistance in the Week 4 note.

\subsubsection{Modeling in AnB}
\begin{lstlisting}
Protocol: OpenAuth_W4_KeyLookup

Types:
  Agent A,I,X;
  Function pk,fKeyReq,fKeyResp;

Actions:
  A -> I: fKeyReq(A,X)
  I -> A: {fKeyResp(X,pk(X))}inv(pk(I))

Goals:
  A authenticates I on fKeyResp(X,pk(X));
\end{lstlisting}

\newpage
\section{TLS Channels, Guessable Password, and Final Status - Mirtha}

\subsection{Outline of the design}

\subsubsection{Week 5: pseudonymous channels}
The Week 5 variant models TLS-like communication through pseudonymous channels.
This abstracts transport-level cryptography while keeping authorization logic explicit.

\subsubsection{Week 6: guessable password checks}
Week 6 evaluates password guessing resistance in three steps:
\begin{itemize}
  \item baseline guessable model (\texttt{guesswhat} attack found),
  \item hardened password-proof model,
  \item focused guessable-only check (bounded no-attack result).
\end{itemize}

\subsection{AnB formalization}

\subsubsection{Design decisions and development}
Week 5 shows channel abstraction alone does not remove logic-level vulnerabilities.
Week 6 introduces dedicated guessing analysis and isolates the required special-task goal in a focused model.

\subsubsection{Modeling in AnB}
\begin{lstlisting}
Protocol: OpenAuth_W6_GuessablePassword_OnlyCheck

Types:
  Agent A,B,I,p;
  Number F,T,R,K,Grant,Photos;
  Function pk,pw,h,fGrant,fData;

Actions:
  p -> A: p,B,F,T,R
  A -> I: {A,p,B,F,T,R,K,h(pw(A,I),R,K)}pk(I)
  I -> A: {fGrant(A,p,B,F,T,Grant)}inv(pk(I))
  A -> p: {fGrant(A,p,B,F,T,Grant)}inv(pk(I))
  p -> B: {fGrant(A,p,B,F,T,Grant)}inv(pk(I)),F
  B -> p: {fData(Photos)}pk(p)

Goals:
  pw(A,I) guessable secret between A,I;
\end{lstlisting}

\newpage
\section{Development and Special Tasks Summary}

\begin{center}
\begin{tabular}{ll}
\toprule
\textbf{Week} & \textbf{Special Task Outcome} \\
\midrule
2 & Static Dolev--Yao analysis completed (baseline weakness identified). \\
3 & Lazy intruder executability of role $P$ completed (lecture-style derivation). \\
4 & Type-flaw resistant format strategy + separate key lookup protocol. \\
5 & Pseudonymous channel variant implemented and analyzed. \\
6 & Guessable-password progression completed; focused bounded check has no guess attack. \\
\bottomrule
\end{tabular}
\end{center}

\section*{AnB files delivered}
\begin{itemize}
  \item \texttt{Week2\_Problem2\_First\_AnB\_Version.anb}
  \item \texttt{Week3\_Problem1\_AnB\_With\_IdentityProvider.anb}
  \item \texttt{Week3\_Problem3\_Refined\_HonestProvider.anb}
  \item \texttt{Week4\_Problem1\_MainProtocol\_Formats.anb}
  \item \texttt{Week4\_Problem2\_AnB\_KeyLookup\_Protocol.anb}
  \item \texttt{Week5\_Problem1\_PseudonymousChannels.anb}
  \item \texttt{Week6\_Problem1\_GuessablePassword\_Check.anb}
  \item \texttt{Week6\_Problem2\_GuessablePassword\_Hardened.anb}
  \item \texttt{Week6\_Problem3\_GuessablePassword\_OnlyCheck.anb}
\end{itemize}

\section*{References}
\begin{itemize}
  \item Sebastian M\"odersheim, Protocol Security Verification Tutorial, DTU Compute.
  \item OFMC and AnB documentation.
  \item 02244 course slides and assignment handout.
\end{itemize}

\end{document}
